% Propietario conservado: /Users/ruben/Documents/New project/output/EXTREMA_MEDIA_RAZON_RECIPROCIDAD_ES_EN_20260918/ARTICULO/ES/source/sections/orientaciones_reticulares.tex
\subsection{Full-support orientations and lattice equivariance}
\label{x_reciprocidad:paper:24-orientaciones}

The weight enumerator of \eqref{x_reciprocidad:exc:enumerador} contains
\(24y^{12}\). Therefore, the set
\[
 C_W^\times:=
 \{u\in C_W:u_b\in\{1,2\}\text{ for every }b\}
\]
has 24 elements. For each \(u\in C_W^\times\), we define
\[
 P_b(u)=
 \begin{cases}
  \mathsf R,&u_b=1,\\
  I,&u_b=2,
 \end{cases}
 \qquad
 g_u:=\bigoplus_{b=0}^{11}P_b(u)\mathsf C P_b(u)^{-1}.
\]

\begin{proposition}[Robustness of Full Support Orientations]
\label{x_reciprocidad:paper:robustez-orientaciones}
For every \(u\in C_W^\times\),
\[
 \frac{g_uv-v}{3},
 \quad
 \frac{g_u^{-1}v-v}{3}
 \in N_0.
\]
In particular, each \(g_u\) preserves the neighbor constructed with the
transported orientation.
\end{proposition}

\begin{proof}
In simple root coordinates,
\[
 \mathsf C\rho-\rho=(-2,-1),
 \qquad
 \mathsf C^{-1}\rho-\rho=(-1,-2).
\]
The discriminant classes of the two displacements are \(u\) and \(-u\).
For the radial coefficient \(a_b\in\{4,1\}\),
\[
 \left\langle
 \frac{(\mathsf C^{\pm1}-I)a_b\rho}{3},a_b\rho
 \right\rangle=-a_b^2.
\]
The sum over the twelve factors again gives \(-27\). The proof of
\cref{x_reciprocidad:km:estabilidad-vecino} is applied component by
component.
\end{proof}

To type the recalibration, let
\(h\in\operatorname{Aut}(C_W)\) be monomial: a permutation \(\sigma\) of
the twelve coordinates followed by multipliers
\(\epsilon_b\in\{1,2\}\). We define its lattice lift
\begin{equation}
 \widetilde h
 :=P_\sigma\circ
 \bigoplus_{b=0}^{11}\mathsf R^{\nu_b},
 \qquad
 \nu_b=
 \begin{cases}
 0,&\epsilon_b=1,\\
 1,&\epsilon_b=2.
 \end{cases}
 \label{x_reciprocidad:paper:elevacion-monomial}
\end{equation}
Its action on the discriminant is exactly \(h\). If the change simultaneously
transports the code, flag, incidence origin, vector \(v\), neighbor, and
frames, then
\begin{equation}
 g_{hu}=\widetilde h\,g_u\,\widetilde h^{-1}.
 \label{x_reciprocidad:paper:naturalidad-isometrias}
\end{equation}
The ternary matrix \(h\) does not act without a lift on \(A_2^{12}\) or on
the Leech lattice.

